\section{What this earns beyond an absolute-forecast pipeline}
The consequence is a decision-specific acquisition theorem. Under its stated structural class, an intervention can be selected with a relative first-event mean guarantee while the absolute lifetime remains statistically unresolved. The missing scale is not approximated and then forgotten: it cancels from the exact comparison. The boundary branch also avoids constructing any experimental Markov chain. Its finite local binary responses are mapped directly into a decision interval.

\subsection{The strongest applicable ordinary methods}
Green identities, Dirichlet and Thomson principles, capacity-based rare-event analysis, and exponential rare-hit approximations are established \cite{aldousfill,bovier,aldousbrown}. The boundary-flux method is an especially strong ordinary comparator in the uniformly reliable large-$N$ family. It is the preferred branch here. It would be misleading to compare only with brute-force waiting for a catastrophe or with reconstructing $2^N$ states.

Loper, Zhou and Geman already reduce suitable rare first-passage probabilities to target-neighborhood capacities and local simulations, with explicit flatness errors \cite{loper}. Their paper also explains why those capacities do not recover an overall first-passage time scale. The present operation exploits a common invariant law so that this missing scale cancels in an intervention comparison, proves a finite Hamming-neighborhood correction from local bounds, and constructs the measurements using unknown operating local cells. We do not claim that local capacity methods or nuisance cancellation are new ideas in themselves.

For independent identical components an ordinary birth--death or product-integral solver is computationally superior to a general rare-path engine. For known numerical coefficients it can give an accurate absolute mean quickly. The lower bound below concerns acquiring its unknown coefficients sufficiently well: granting the best such exact solver does not supply the needed information. It compares with every absolute-mean acquisition procedure in the same local access class, not just a naive uniform confidence-box estimator.

The accepted marked sampler and CDF certificate remain stronger output capabilities. They support general event predicates and other losses under their scopes. This new operation saves work by committing to one objective, one common-invariant intervention family, and a singleton event; it does not make those broader forecasts unnecessary for questions they actually answer. No universal superiority over every specialist direct-decision algorithm is claimed. The explicit boundary branch is one such specialist method, now joined to finite acquisition, errors and deployment.

\section{A same-access separation between a decision and an absolute forecast}
\label{sec:sameaccesslower}

This lower bound concerns acquired information, not the running time of a
solver given exact weights.  It grants the forecaster strictly more
observations than the finite-hold decision needs.  In particular it grants
perfect clocks, arbitrary calibrated additive fields, and complete continuous
observation of each released local cell.  The obstruction therefore does not
come from a slow unknown clock or from requiring an inefficient parameter
estimator.

\paragraph{The experiment class.}
There are $n$ independent binary components, with unknown logits
$\theta_i$, equilibrium failure probabilities $q_i=\sigma(\theta_i)$,
and \emph{known unit refresh clocks}.  At applied field $u$, a local cell
has rates
\[
  0\longrightarrow1:\ \sigma(\theta_i+u),\qquad
  1\longrightarrow0:\ 1-\sigma(\theta_i+u).
\]
An experiment may adaptively choose a component, an arbitrary field, a
known initial bit, a release duration, and its stopping rule.  It may observe
the complete resulting local path.  Several cells may operate in parallel;
their free durations are added.  A supplied fresh equilibrium bit at a
chosen field is also permitted and counted as one local observation.  Pins,
controller randomization, and computations convey no additional unknown
parameter by themselves.  An equilibrium bank of $n$ bits counts as $n$
local observations; parallelism does not turn its $n$ sensing channels into
one unit of acquired evidence.  These are the declared local state/clock
experiments, not direct measurements of unknown Hamiltonian coefficients.

Write
\[
  W=B+r^*\sum_{\text{released cells}} \text{free duration},
\]
where $B$ is the number of supplied fresh equilibrium bits and $r^*$ is
the known maximum available refresh clock (one in the unit-clock model).
Allowing additional known clock modules up to a fixed $r^*$ only changes
this normalization: their rates are the same under the two hypotheses.
This relaxed
bill omits pin, detector, field, memory, and computational costs, so a lower
bound on it remains a lower bound when those services are charged.

\paragraph{Theorem (forecast information lower bound).}
Let $n\ge4$, $0<\epsilon\le1/8$, and $0<\delta<1/2$.
Suppose an adaptive procedure returns an estimate $\widehat M$ satisfying
\[
 \Pr_\theta\{|\log\widehat M-\log M(\theta)|\le\epsilon\}\ge1-\delta
\]
for every allowed model, where $M$ is the mean first visit to
$b=(1,\ldots,1)$ from the equilibrium law conditioned on $b^c$.
Then, on at least one allowed model,
\begin{equation}
  \mathbb E W\ \ge
  \frac{n^2}{512e\,\epsilon^2}
       \operatorname{kl}(1-\delta,\delta).
  \label{eq:forecastlower}
\end{equation}
With only fresh equilibrium-bit experiments the stronger constant
$n^2\operatorname{kl}(1-\delta,\delta)/(128\epsilon^2)$ holds.
For fixed $\delta<1/2$, this is an $\Omega(n^2/\epsilon^2)$ local-evidence
requirement.  It holds even though the operating gap equals one, the clocks
are known, and every coordinate can be actively balanced.

\paragraph{Proof.}
First consider equal logits $\theta_i=\theta$, and put $q=\sigma(\theta)$,
$p=q^n$.  Independence and the two-state semigroup give
\[
 P_t(b,b)=\{q+(1-q)e^{-t}\}^n.
\]
The classical fundamental-matrix identity, also derived in the return
certificate below, gives
\begin{align}
 S_n(q)&=\int_0^\infty[P_t(b,b)-p]\,dt
       =\sum_{k=1}^n {n\choose k}
                    \frac{q^{n-k}(1-q)^k}{k},\\
 M_n(q)&=\frac{S_n(q)}{p(1-p)}.
 \label{eq:productmeanlower}
\end{align}
Let $K\sim\operatorname{Bin}(n,1-q)$.  Then
$S_n=\mathbb E[K^{-1};K>0]$.  Since
$1/k\le2/(k+1)$ for $k\ge1$,
\[
 S_n(q)\le\frac{2}{(n+1)(1-q)}.
\]
Direct differentiation of the finite sum gives
\begin{equation}
 -\frac{d}{d\theta}\log M_n(\sigma(\theta))
 =\frac{1-p}{S_n(q)}-\frac{np(1-q)}{1-p}.
 \label{eq:meanlogderivative}
\end{equation}
For $n\ge4$ and $q\le1/2$, the right side is at least
\[
 \frac{15(n+1)}{64}-\frac n{15}\ \ge\ \frac n8.
\]
Take $\theta^{(0)}_i=-\log3$ and
$\theta^{(1)}_i=-\log3+d$, where $d=32\epsilon/n\le1$.
The probabilities along this segment stay below $1/2$.
Equation~\eqref{eq:meanlogderivative} therefore yields
\[
 \log M(\theta^{(0)})-\log M(\theta^{(1)})\ge4\epsilon.
\]
A valid estimate distinguishes these two hypotheses with both error
probabilities at most $\delta$.

For one equilibrium-bit observation at any adaptively selected field,
the Bernoulli log-partition function has second derivative at most $1/4$.
Consequently its conditional divergence between the two hypotheses is at
most $d^2/8$, independently of the field magnitude.  For a released cell,
write $\lambda_0(x,u),\lambda_1(x,u)$ for the applicable exit rates.
Because the derivative of each logistic log rate has magnitude at most one,
\[
 \left|\log\frac{\lambda_1(x,u)}{\lambda_0(x,u)}\right|\le d.
\]
The continuous-path likelihood identity gives conditional divergence
\[
 \mathbb E_0\int
 \left[\lambda_0\log\frac{\lambda_0}{\lambda_1}
             -\lambda_0+\lambda_1\right]dt
 \le\frac{ed^2}{2}\,\mathbb E_0\int dt,
\]
since $\lambda_0\le1$ in the unit-clock case (or at most $r^*$ with
additional clock modules) and
$e^v-1-v\le e d^2/2$ for $|v|\le d\le1$.
Adaptive fields, adaptive stopping, and repeated use of a cell obey the
same conditional likelihood chain rule.  Truncating an unbounded stopping
rule and then taking a limit gives
\[
 \operatorname{KL}(\mathsf P_0^{\rm data}\Vert\mathsf P_1^{\rm data})
 \le\frac{ed^2}{2}\,\mathbb E_0 W.
\]
Data processing through the final binary test lower-bounds this divergence
by $\operatorname{kl}(1-\delta,\delta)$, proving
\eqref{eq:forecastlower}.  The equilibrium-only bound uses $d^2/8$ in
place of $ed^2/2$.  Finite instrument errors cannot increase the information
available beyond this deliberately idealized observation class.
\hfill$\square$

\paragraph{Heterogeneous and nontrivial decision instances.}
The hard pair need not consist of identical components.  Let initial
$q_i$ belong, for example, to $\{1/4,1/3\}$ in any unknown arrangement.
Apply the same additive logit shift $t$ to every component.
If $K$ is the resulting Poisson-binomial number of working components,
$S=\mathbb E[K^{-1};K>0]$ and $p=\prod_iq_i$ satisfy
\[
 \frac{dS}{dt}=\Lambda S-(1-p),\qquad
 \frac{d\log p}{dt}=\Lambda,
 \qquad \Lambda=\sum_i(1-q_i).
\]
Thus
\[
 -\frac{d\log M}{dt}=\frac{1-p}{S}-\frac{p\Lambda}{1-p}.
\]
As long as all $q_i\le1/2$, stochastic domination by
$\operatorname{Bin}(n,1/2)$ gives the same lower bound $n/8$.
For $n\ge8$, $\epsilon\le1/8$, the shift $32\epsilon/n$ keeps the
displayed heterogeneous family below $1/2$.  The information bound is
unchanged.  These models can have a nontrivial optimal assignment of unequal
clock modules, while the unknown common energy shift leaves their optimal
ordering unchanged.  The lower bound is therefore not an artifact of
uncertainty about a globally rescaled time unit.

\paragraph{Matching order for prediction in this product family.}
The $n^2$ dependence is attainable up to confidence and finite-service
logarithms; it is not an inverse-$p$ assertion.  For $q_i$ in a fixed compact
subinterval of $(0,1)$, sample each local equilibrium probability $m$ times,
clip its empirical frequency to a slightly larger compact subinterval,
and add the $n$ log frequencies.  The bias is $O(n/m)$ and bounded
differences give fluctuation
$O(\sqrt{n\log(1/\delta)/m})$.  Taking
$m=O(n\epsilon^{-2}\log(1/\delta))$ estimates $\log p$ to $O(\epsilon)$
using $O(n^2\epsilon^{-2}\log(1/\delta))$ local observations.  The
nonrare return prefactor $S$ is acquired separately by the finite return
operation.  Every ideal equilibrium read can be realized by a finite hold
at the known unit clock, with conditional bias below its allotted error;
this introduces an explicit logarithmic dwell factor, not $1/p$.

\paragraph{Meaning of the separation.}
At fixed margins and fixed decision accuracy, the direct acquired
intervention constructions use $\widetilde O(n)$ local evidence, while
\emph{every} uniformly accurate absolute-mean forecaster in the same local
access class has the $\Omega(n^2)$ requirement above.  A method is free to
use the same nuisance cancellation and become a direct decision method;
the theorem does not claim a lower bound against such a method.  It shows
that requiring a calibrated absolute forecast before deciding is an
additional statistical task.  It is not a claim that ordinary simulation
is the strongest baseline, nor that all rare-event decisions are cheap.

The proof uses classical likelihood/change-of-measure tools, as in
fixed-confidence experimental identification.  The established
fundamental-matrix and bandit results are machinery; no priority claim is
made for those identities.  The derived separation is for the explicit
rare first-event intervention family and the paid local experiment class
above.


For direct comparison with the numerical family's $\bar q=1/4$, the same proof can use $\theta^{(0)}=-\log7$ and $N\ge8$. Its shifted probabilities remain below $1/4$ for $\epsilon\le1/8$, while the derivative and information bounds used above still hold. Thus the quadratic forecast lower bound also holds inside a fixed uniformly reliable subclass with the same kinds of coarse certificates, known fast clocks and safe preparation. The output bit length, input description and parallel hardware costs are not removed by either method.

\section{Executed growing-family decision and its limits}
\label{sec:execution}
The root program implements only the data flow required by the decision. Its source simulator is hidden behind an actual-context local-hold interface. It supplies binary outcomes, and the decision receives public coarse bounds, menu labels and $N$. It never reads the numerical energies, clocks, stationary mass or absolute lifetime.

The family has nonuniform fields $h_i=-2+0.25\sin(0.73i)$ and ring interactions $J_{i,i+1}=0.12\cos(0.41i+0.3)$. The three clock menus cross between contexts. One has clocks in $[0.25,0.30]$ when both neighbors have failed and clocks in $[0.85,1]$ elsewhere; a second reverses those assignments. A third has a different heterogeneous context response in $[0.5,0.65]$. Consequently there is no coordinatewise generator order deciding the menu. All share the same Gibbs law. Valid loose certificates are $\bar q=1/4$, $r_-=0.2$, $r_+=1$, $\alpha=0.7$ and $\gamma=0.14$.

Each run uses $4096$ binary local holds for each of three actions. The nominal hold is $h=1/4$, its exogenous deadline varies with prior history by at most $0.0025$ time units, and the simulator adds read corruption with probability $0.0005$. The certificate permits combined conditional bias $0.001$. It uses an exact-integer Chernoff inversion and exact-rational structural inequalities. A tail budget $1/1200$ per item, tail and size gives simultaneous statistical confidence at least $0.98$ across all four runs in the ideal random experiment. The executed seeded PRNG record is a reproducible diagnostic, not a certified physical random source.

\begin{center}
\begin{tabular}{rrrrr}
\toprule
$N$&Boundary $k$&Upper $q_0S$&Mean ratio vs fast menu&Mean ratio vs mixed menu\\
\midrule
64&14&1.053914&1.735710&1.045747\\
256&15&1.009896&1.860830&1.055450\\
1024&17&1.002926&1.904448&1.257433\\
4096&18&1.002223&1.717734&1.080755\\
\bottomrule
\end{tabular}
\end{center}
The last two columns are conservative certificate lower bounds on the selected action's true equilibrium-safe mean divided by the competitor's mean, before the separately specified deployment/preparation allowance. They are not estimates of the actual improvements. Every run certifies the same menu as best. The state spaces range from $2^{64}$ to $2^{4096}$; the proof, rather than enumeration, governs the growing family.

Every size uses $12\,288$ local holds and $3061.76$ total free-cell time units. No full-system failed preparation and no virtual trajectory were used. Additional pin, neighbor-context, read, timer, field, setup, controller and installation services are not included as zero. The experiment assumes the validated coarse certificates; their full finite acquisition is constructed above and its conservative bill is evaluated separately below.

\subsection{Composition with the stronger mean-only initialization}
The file \texttt{deployment\_composition.json} applies the all-safe-start
theorem to the same recorded counts using exact rational arithmetic.
The preferred branch runs the actual installation for a time in
$[2,2.001]$ with no inspection or rejection. The graphical comparison
gives initial-failure mass at most
$(13/16)^N+(N/4000)(13/16)^{N-1}$, including the entire clock window.
It is below $1.727\,10^{-6}$ at $N=64$ and decreases exponentially.
The actual selected-to-best mean is then at least $0.9904350$ at
$N=64$; the chosen menu is certified best at the other three sizes.
Whole-system free evolution is capped at $2.001$ model-time units,
and summed component exposure at $2.001N$. This evaluated interface
uses the proved uniform rate bounds and repairable commissioning;
it was not executed on physical hardware and does not assert a
certainly safe initial specimen or equilibrium-law preparation.

A separately evaluated meter contract has captured-bit error at most
$1/4000$, capture-to-mission delay at most $0.001$, and actual warm-up
in $[2,2.001]$ model-time units. It gives working probability at least
$3/16$, acceptance at least $11997/64000$, and conditional initial
failure at most $32/11997<0.002668$. Thirty-four attempts bound the
preparation-abort probability by $0.001$. Their whole-system warm-up
and latency cap is $68.068$ time units, costing $68.068N$ component-time
before the other physical services. This is an evaluated interface,
not executed physical preparation.

The selected-to-best actual successful-preparation mean is certified
at least $0.9877949$ at $N=64$. At $N=256,1024,4096$, the selected
menu is still certified best after these additional allowances. If
preparation abort receives zero lifetime, the corresponding selected-to-best
guarantees are $0.9868071$ at $N=64$ and at least $0.999$ for the other
sizes. Thus the equilibrium-safe table is not silently reused as an
unchanged physical-start certificate. No full initial law is claimed.

\subsection{Separate grouped-wait execution}
The completed \texttt{next6\_rate\_assay} program executes finite
censoring, clipping, persistent clock offsets and exact rational output
arithmetic on heterogeneous local contexts at $N=64,4096,262144$.
It uses $49\,152$ groups per menu at relative tolerance $1/8$, for six
menu assays with simultaneous confidence at least $0.95$ conditional
on the coarse bounds. All three pairs have disjoint confidence intervals.
The six executions used $294\,912$ groups, $1\,179\,647$ pin--release
observations and $1\,597\,822.63$ active model-time units. The one
censored group and its discarded exposure are included. Observed relative
errors below $0.54\%$ do not replace the $12.5\%$ certificate.
This is a local-interface execution, distinct from the interacting-ring
global decision and from a physical apparatus experiment.

\subsection{A full bill can still be large}
The file \texttt{calibration\_costs.json} evaluates a simple sufficient coarse acquisition with hold $5$, timing error at most $0.05$, per-indicator statistical radius $0.01$, service bias $0.0005$, and total calibration error probability $0.005$. It includes all $4N$ contexts for each of three clock menus and both starting bits. This schedule was evaluated, not experimentally executed.
\begin{center}
\begin{tabular}{rrrr}
\toprule
$N$&Coarse pins/reads&Coarse free exposure&System prep time cap\\
\midrule
64&$1.02364\,10^8$&$5.16936\,10^8$&$1.89661\,10^4$\\
256&$4.52039\,10^8$&$2.28280\,10^9$&$7.46155\,10^4$\\
1024&$1.97849\,10^9$&$9.99138\,10^9$&$2.97213\,10^5$\\
4096&$8.59541\,10^9$&$4.34068\,10^{10}$&$1.18760\,10^6$\\
\bottomrule
\end{tabular}
\end{center}
Exposure is in model time units, not automatically seconds. The preparation column is whole-system free-evolution time, using a pointwise density-log allowance $0.01$, a certified conditional floor $q_-=0.025$, and eleven capped attempts. Its summed component exposure is $N$ times that column; inspection, installation and field ramping are additional. If one time unit were one second, the first coarse exposure alone would exceed sixteen serial cell-years. This is a deliberately exposed sufficient cost, not a lower bound and not a claimed useful device. The inexpensive postcertificate decision must not be presented as the entire acquisition bill. The theorem's established advantage is asymptotic information scaling and removal of unnecessary forecast requirements; a practical regime still depends on the actual source and services.

\subsection{Useful unresolved attempts retained}
The first numerical certificate used the larger class of timers allowed to stop in response to the current focal path. Its weaker coefficient left three of four menu comparisons unresolved at the chosen sample count. That record is preserved in \texttt{stronger\_timer\_contract\_unresolved.json}. The executed source uses exogenous deadlines; applying the appropriate, separately proved bound certifies all four using the same underlying outcome counts. The final coarse envelopes were also loosened to have explicit acquisition slack. No outcome was discarded and no numerical failure was renamed a guarantee.

The initial candidate architecture used return probes for every decision. The ordinary boundary-flux comparator made that unnecessary in the certified regime and replaced it. Earlier attempts to use only generator order or convexity were retained as limited mechanisms: they do not solve arbitrary preventive allocation. The distinctions between a mean decision, a CDF, a marked record and a physical intervention are maintained throughout.

\section{Independent completed consequences}
\subsection{All factorially many assignments in a narrower physical model}
The independent file \texttt{next6\_decision\_theory/sorting\_intervention.tex} proves a global assignment rule for independent components with transferable reciprocal clock modules. Installing module $j$ at component $i$ must produce actual rates $r_jq_i$ and $r_j(1-q_i)$. Pairing larger $q_i$ with larger $r_j$ maximizes the equilibrium-safe mean to the all-failed state among all $N!$ assignments. The opposite order minimizes it. The exchange identity
\[
 [q_i+(1-q_i)e^{-at}][q_j+(1-q_j)e^{-bt}]
 -[q_i+(1-q_i)e^{-bt}][q_j+(1-q_j)e^{-at}]
 =(q_i-q_j)(e^{-bt}-e^{-at})
\]
proves the pointwise return comparison. This exact ordering survives an arbitrary independent ergodic spectator process with a specified singleton target; its preparation costs remain.

A dimension-independent bound on the logarithmic sensitivity of the nonrare return prefactor means constant, decision-level probability accuracy suffices. The completed acquisition uses $O(N\epsilon^{-2}\log(N/\delta))$ local reads at fixed margins, plus priced holds, timers and installation, followed by $O(N\log N)$ sorting. It does not need $1/N$ parameter accuracy or an absolute lifetime. Classical rearrangement and reliability-allocation precedents are credited; no general clock-module transport is inferred from biased calibration.

\subsection{Continuous design for desired completion}
The separate \texttt{next6\_direct\_decision} construction minimizes a desired rare completion time over a rational polytope of certified reciprocal kinetic multipliers. Its normalized objective $pR M(c)$ is convex with rarity-free bounds and curvature. A finite scalar oracle and inexact Frank--Wolfe optimize the continuum without a control grid or global matrix. This is completion acceleration, not maximizing a convex prevention objective. Bounded multipliers improve actual completion by at most their bounded factor.

Its evaluated elementary sufficient schedule at $N=128$, accuracy $0.1$, uses about $4.12\,10^8$ scalar calls and $10^{61.65}$ potential updates. These were not executed and are not lower bounds. Product-reference executions have a stronger ordinary solver and gains only about $1.10$--$1.23$. This useful global-design deduction was therefore not selected as the main capability.

\subsection{Quantum export without Hamiltonian identification}
The independent \texttt{next6\_physical/CONTROL\_FOR\_CALIBRATION.md} adds calibrated local source rotations to the accepted pure coherent Weyl ports. Covariant source/output frame correction and fresh Pauli averaging compare the full retained dilation, not only the reduced source. For unknown predetermined $H(t)$ with $\|H(t)\|\le\Lambda$, its finite export error obeys
\[
 2Ne^{-rT}+\frac{2\Lambda^2T^2}{B}
       +2(\Lambda+Nr)W_{\rm lost}+\epsilon_{\rm interface}.
\]
It works for every $r>0$ without Hamiltonian-learning exposures, with finite smooth pulses and charged lost output intervals. The random key is corrected and excluded from the returned target; retaining it does not have the same quadratic guarantee. Pure ports, large calibrated coherent receivers, protected memory, source strength and slew remain substantive. At $N=64$, $\Lambda=N$, $r=1$, $\epsilon=0.1$, a sufficient recipe uses about $23.9$ million blocks and peak per-site control norm $7.69\,10^{11}$ before receiver costs. Established randomized decoupling is credited. This is an additional-access resource trade, not same-access dominance over output-only acquisition or an economical receiver.

\section{Preservation and scientific status}
All 229 files of the preceding acquired-first-event package are preserved byte for byte, including its complete fitted-model and experiment-fed laws, operating and static coin constructions, confidence/preparation bills, independent results and nested earlier sources. The new archive's preservation receipt checks the original workspace, copied originals and ZIP members against the initial SHA-256 baseline. Historical failed or incomplete attempts remain historical files; the newest governing reports specify which interfaces were subsequently completed.

The result earned here is a complete conditional theorem, finite acquisition/decision/deployment construction, explicit same-access forecast separation, uniform safe-start mean comparison, and growing-family numerical execution. It is not evidence that a natural system satisfies the model, a prototype, a universal prevention optimizer, or demonstrated civilization-scale impact. The scientific principle is precise: when interventions share an invariant rare-event scale, the local kinetic information needed to decide can be much smaller than the information needed to forecast that scale. In the certified boundary regime even the initial full law can be discarded for this objective, because the theorem controls the mean uniformly over safe states.
